% Generated draft: requires mathematical and source review.
\documentclass[11pt]{article}
\usepackage[margin=25mm]{geometry}
\usepackage{fontspec}
\setmainfont{Latin Modern Roman}
\usepackage{amsmath,amssymb,mathtools,unicode-math}
\setmathfont{Latin Modern Math}
\usepackage{enumitem,needspace,xurl,hyperref}
\setlength{\parindent}{0pt}
\setlength{\parskip}{5pt}
\setlength{\emergencystretch}{3em}
\newcommand{\N}{\mathbb{N}}
\newcommand{\Z}{\mathbb{Z}}
\newcommand{\Q}{\mathbb{Q}}
\newcommand{\R}{\mathbb{R}}
\newcommand{\C}{\mathbb{C}}
\newcommand{\F}{\mathbb{F}}
\newcommand{\A}{\mathbb{A}}
\newcommand{\PP}{\mathbb{P}}
\newcommand{\E}{\mathbb{E}}
\newcommand{\eps}{\varepsilon}
\newcommand{\dd}{\,\mathrm{d}}
\begin{document}
\section*{UnsolvedMath ID: 8800034}
\textbf{Source problem code:} catalogue entry
\subsection{Definitions and mathematical statement}
Let $q$ be a prime power and let $E/\F_q$ be an elliptic curve. Its rational-point group has the form
\[
E(\F_q)\cong\Z/m\Z\times\Z/n\Z,\qquad m\mid n,
\]
with $m\mid q-1$ and $|mn-(q+1)|\leq2\sqrt q$. These are necessary structural constraints, not an assertion that every pair satisfying them is realized.

For example, under uniform sampling of $\F_q$-isomorphism classes one can define
\[
C_q(m,n)=\#\{[E]_{\F_q}:E(\F_q)\cong\Z/m\Z\times\Z/n\Z\},\qquad
P_q(m,n)=\frac{C_q(m,n)}{\sum_{a,b}C_q(a,b)}.
\]
The source asks to study this type of distribution and identify its correct nonuniform law. It does not prescribe uniform isomorphism-class sampling: weighting each curve by $1/|\operatorname{Aut}_{\F_q}E|$, sampling equation coefficients, or conditioning on ordinary/supersingular curves gives different measures. The display is one explicit candidate model.

A limiting law also needs an asymptotic regime and a normalization of the changing groups: $q$ through primes, all prime powers, or fixed-characteristic extensions are different questions. Fixed-field exact counts, asymptotic frequencies, and pointwise uniform estimates are distinct targets. No such regime, error term or conjectural density is supplied, so the request is a distribution programme rather than one specified limiting theorem.
\subsection{Short English statement}
Determine how elliptic-curve rational-point group types are distributed, using a specified sampling measure and finite-field asymptotic regime. Exact counts, cyclic frequencies and a limiting nonuniform law are separate questions.
\subsection{Sources}
\begin{itemize}
\item Local UnsolvedMath registry, permanent ID 8800034, exact source code catalogue entry, ``Distribution of elliptic-curve group structures''; attributed to Igor Shparlinski, no proposed year supplied. Registry status partially solved, classification PARTIAL-PROGRESS.
\item \emph{IPAM Workshop on Number Theory and Cryptography---Open Problems} (2006), attachment 6224, PDF page 9, \url{https://www.ipam.ucla.edu/programs/workshops/workshop-i-number-theory-and-cryptography-open-problems/?tab=schedule}. Supplied pdf-and-legacy-ppt visual audit accessed July 29, 2026; provenance records 17 PDFs/two PPT decks inspected with repeated builds excluded. No PDF/PPT was accessed here.
\item The summary names Lenstra (1987) for an ``essentially uniform'' Hasse-interval order statement and Vladut (1999) for an at-least-75-percent cyclic statement. No sampling, limiting regime or exact hypotheses are supplied, so neither phrase is adopted as an unqualified fixed-$q$ probability theorem.
\item ``On group structures realized by elliptic curves over a finite field'', \texttt{arXiv:1003.3000}, is supplied for exact structure counts and realized-type counts, without authors. The summary also names David--Smith for conditional asymptotics using primes in short arithmetic progressions and unspecified \emph{Canadian Journal of Mathematics} follow-up bounds, without full references.
\item The supplied assessment reports substantial progress but no settled general nonuniform limiting law. Provenance flags current primary-literature status and redistribution terms NEEDS\_REVIEW.
\end{itemize}
\end{document}
